\section{Geometric consequences and examples}\label{cons:section}\label{sec:consequences}

The criterion separates two different questions about a spherical set:
whether its ordinary quadratic equations admit the required identity
block, and whether that solution lifts to the tensor ring. We first give
positive examples, then exhibit two failures of the lift.

\begin{corollary}\label{cons:spherical}
Every finite Euclidean Ramsey set is spherical.
\end{corollary}
\begin{proof}
A singleton is spherical. Otherwise use the notation and a matrix $P$
from Theorem~\ref{thm:classification}, and put $H=m_F(P)$, entrywise.
Multiplying the evaluation equations gives
\[
 \|a_i\|^2+\ell\cdot a_i+c=0,
 \qquad c=H_{00},\qquad
 \ell_\alpha=H_{0\alpha}+H_{\alpha0}\quad(1\le\alpha\le d).
\]
Thus all points lie on the sphere with center $-\ell/2$ and squared
radius $\|\ell\|^2/4-c$. The radius is positive because the set has
at least two distinct points. No symmetry of the constant row and
column of $H$ is needed.
\end{proof}

A finite Euclidean set is \emph{transitive} if its distance-preserving
permutations act transitively on its points.

\begin{corollary}\label{cons:transitive}
Every nonempty subset of a finite transitive Euclidean set is Ramsey.
\end{corollary}
\begin{proof}
Singletons are immediate. Let $A=\{a_1,\ldots,a_s\}\subset\R^d$,
with $s\ge2$, affinely span $\R^d$, and suppose an isometric
embedding maps $a_i$ to $y_i$ in a finite transitive set
$Y\subset\R^n$. Write $H$ for the finite group of distance-preserving
permutations of $Y$. For each $h\in H$, the map $a_i\mapsto h(y_i)$
extends to an affine isometry. Let $w_{hj}\in W=\R^{1+d}$ be its
$j$th affine coordinate row, so that $e_i(w_{hj})=(h(y_i))_j$.
The gradient Gram matrix of these $n$ rows is $I_d$.

Let $c_{hj}=((h(y_1))_j,0)\in W$, a constant affine row. The
symmetric rational tensor
\[
 T=\frac1{|H|}\sum_{h\in H}\sum_{j=1}^n
       \bigl(w_{hj}\otimes w_{hj}-c_{hj}\otimes c_{hj}\bigr)
\]
has $G(T)=I_d$. For each $i$, the multiset $(h(y_i))_{h\in H}$
contains every point of $Y$ exactly $|H|/|Y|$ times. Consequently
\[
 \sum_{h,j}(h(y_i))_j\otimes(h(y_i))_j
 =\sum_{h,j}(h(y_1))_j\otimes(h(y_1))_j,
\]
and $(e_i\otimes e_i)(T)=0$. Proposition~\ref{alg:equivalence}
and Theorem~\ref{thm:classification} now give the Ramsey property.
\end{proof}

Quadratic interpolation gives another source of tensor certificates.
The row-independence hypothesis below says that quadratic polynomials
can prescribe arbitrary values on the configuration.

\begin{proposition}\label{cons:quadratic}
Let $A=\{a_1,\ldots,a_s\}\subset\R^d$ be spherical and affinely
span $\R^d$. Put $p_i=(1,a_i)^{\mathsf T}$ and let $F$ be its
coordinate field. If the $s$ row vectors
\[
 (p_{i\alpha}p_{i\beta})_{0\le\alpha,\beta\le d}
 \quad(1\le i\le s)
\]
are linearly independent over $F$, then $A$ is Ramsey.
\end{proposition}
\begin{proof}
The singleton case is immediate, so assume $s\ge2$. Write
$N=(d+1)^2$. Sphericity gives a real matrix with spatial block
$I_d$ and quadratic values zero at every $p_i$. These requirements
are a linear system over $F$, so a solution exists over $F$ as well,
by Gaussian elimination. Denote such a matrix, written as a vector in
$F^N$, by $h$.

Let $D\in\operatorname{Mat}_{s\times N}(B)$ be the tensor evaluation
matrix, whose entry in row $i$ and column $(\alpha,\beta)$ is
$p_{i\alpha}\otimes p_{i\beta}$. Lift $h$ to $b\in B^N$ by applying
$x\mapsto x\otimes1$ entrywise. The assumed row independence allows
us to select $s$ columns of $D$ whose square submatrix $M$ has
\[
 \delta=m_F(\det M)\ne0.
\]
Let $\iota:B^s\to B^N$ insert a vector into these selected coordinates,
so that $D\iota=M$. Multiplying by $\det M$ allows us to cancel the
tensor evaluation error $Db$ through the selected columns: the adjugate
identity gives
\[
 z=(\det M)b-\iota\operatorname{adj}(M)Db\in\ker D.
\]
Since $m_F(D)h=0$, multiplication gives $m_F(z)=\delta h$.
Multiplying $z$ by $\delta^{-1}\otimes1$ therefore produces a
solution of the criterion in Theorem~\ref{thm:classification}.
This argument does not require $\det M$ to be a unit of $B$.
\end{proof}

For at most five circle points, products of two line equations supply
the required interpolation. The same observation appears in P\'alv\"olgyi's
analysis of derivation obstructions
\cite[proof of Theorem~A.2]{palvolgyi-heptagon}.

\begin{corollary}\label{cons:quadrilateral}
Every nonempty set of at most five distinct points on a circle is Ramsey.
In particular, every cyclic quadrilateral is Ramsey.
\end{corollary}
\begin{proof}
Work in the plane of the circle and extend the set, if necessary, to
five distinct points on that circle. Let $F$ be the coordinate field of
the resulting five-point set. Fix one of these points $a_i$, partition
the other four into two pairs, and take the product of affine equations for the two lines
through those pairs. No three circle points are collinear, so this
quadratic polynomial vanishes at the other four points and is nonzero
at $a_i$. Its coefficients lie in the coordinate field $F$. Rescaling
the five resulting polynomials shows that quadratic evaluation onto
$F^5$ is surjective. Its five rows are therefore independent, and
Proposition~\ref{cons:quadratic} makes the five-point set Ramsey.
A monochromatic copy of that set contains a copy of the original set.
\end{proof}

\begin{corollary}\label{cons:kites}
For every transcendental $a\in(-1,1)$, the cyclic quadrilateral
\[
 K_a=\bigl\{(-1,0),(1,0),
             (a,\sqrt{1-a^2}),(a,-\sqrt{1-a^2})\bigr\}
\]
is Ramsey and is not subtransitive.
\end{corollary}
\begin{proof}
The four points are distinct and lie on the unit circle, so
Corollary~\ref{cons:quadrilateral} makes $K_a$ Ramsey.
Leader, Russell and Walters proved that $K_a$ does not embed in any
finite transitive Euclidean set \cite[Corollary~2]{lrw-kites}.
\end{proof}

Thus the necessity direction of the Leader--Russell--Walters
subtransitive characterization fails \cite[Conjecture~A]{lrw-transitive}.
The same family also disproves their conjecture that these kites are
not Ramsey \cite[Conjecture~3]{lrw-kites}.

\subsection{Nine points on a circle}

The tensor equations can have full rank even when their multiplied
versions admit a sphere equation. The following determinant calculation
illustrates this distinction. P\'alv\"olgyi's earlier generic-seven-point
claim \cite[Theorem~A.3]{palvolgyi-heptagon} already implies its
non-Ramsey conclusion by taking a subset; the source labels its appendix
as unchecked. We give a direct proof from the tensor criterion.

\begin{proposition}\label{cons:nine}
Let $t_1,\ldots,t_9$ be algebraically independent real numbers over
$\Q$. The nine points
\[
 a_i=\left(\frac{1-t_i^2}{1+t_i^2},
                \frac{2t_i}{1+t_i^2}\right),\qquad 1\le i\le9,
\]
form a spherical set that is not Ramsey.
\end{proposition}
\begin{proof}
The parametrization lies on the unit circle and is injective, since
$t_i=(a_i)_2/(1+(a_i)_1)$. Thus the points are distinct, affinely span
$\R^2$, and have coordinate field $F=\Q(t_1,\ldots,t_9)$.
Algebraic independence identifies this field with a rational function
field. Consequently $B=F\otimes_\Q F$ is the localization of
$\Q[T_1,\ldots,T_9,U_1,\ldots,U_9]$ obtained by inverting all
nonzero polynomials in the $T$ variables alone and all nonzero
polynomials in the $U$ variables alone. In particular, $B$ is a domain
and embeds in $\Q(T_1,\ldots,T_9,U_1,\ldots,U_9)$.

Put
\[
 p(t)=\left(1,\frac{1-t^2}{1+t^2},\frac{2t}{1+t^2}\right)^{\mathsf T}.
\]
Over this fraction field, the nine tensor evaluation equations have
coefficient matrix $D$ with rows
$(p_\alpha(T_i)p_\beta(U_i))_{0\le\alpha,\beta\le2}$.
We claim that $\det D\ne0$. To check the claim, specialize the nine
pairs $(T_i,U_i)$ to the grid $\{0,1,-1\}^2$, with both coordinates
ordered as $0,1,-1$.
All denominators in the displayed matrix remain nonzero. The resulting
matrix is the ordinary Kronecker product $C\otimes C$, where
\[
 C=\begin{pmatrix}1&1&0\\1&0&1\\1&0&-1\end{pmatrix},
 \qquad \det C=2.
\]
Its determinant is $2^6=64$, proving the claim. This evaluates the
specific determinant expression; no specialization of the entire
localized ring $B$ is asserted.

Thus the only solution of $D\,\operatorname{vec}(P)=0$ over the
fraction field is $P=0$. Injectivity of the embedding of $B$ shows
that the same holds over $B$. Such a matrix cannot have multiplied
spatial block $I_2$, so Theorem~\ref{thm:classification} excludes the
Ramsey property.
\end{proof}

\subsection{Three concentric squares}

A second example uses only one transcendental parameter. It also shows
how a derivation can detect the failure of the tensor equations directly
inside the coordinate field. Its weighted cancellation is related to
the derivation obstruction in P\'alv\"olgyi's smaller circular example
\cite[Theorem~1]{palvolgyi-heptagon}; the calculation below is independent
of that proof. The parameter is Liouville's constant
\cite{liouville1851}.

\begin{proposition}\label{cons:twelve}
Set
\[
 t=\sum_{m=1}^{\infty}10^{-m!},\qquad
 u=\frac{1-t^2}{1+t^2},\qquad v=\frac{2t}{1+t^2},
 \qquad Q_0=\{(\pm1,0),(0,\pm1)\},
\]
and define
\[
 R_+=\begin{pmatrix}u&-v\\v&u\end{pmatrix},\qquad
 R_-=\begin{pmatrix}u&v\\-v&u\end{pmatrix},\qquad
 X=R_+Q_0\cup R_-Q_0\cup Q_0.
\]
Then $X$ is a twelve-point subset of the unit circle that is not Ramsey.
\end{proposition}
\begin{proof}
First, $0.11<t<0.12$ and $t$ is transcendental. Indeed, the rational
partial sums $t_m=p_m/q_m$, with $q_m=10^{m!}$, satisfy
\[
 0<t-t_m<2q_m^{-(m+1)}\qquad(m\ge2).
\]
If $t$ had an irreducible polynomial $P\in\Z[T]$ of degree $d$, then
$P(t_m)\ne0$ and $|P(t_m)|\ge q_m^{-d}$. A bound $C\ge1$ for
$|P'|$ on $[0,1]$ would give
$q_m^{-d}<2Cq_m^{-(m+1)}$, a contradiction for large $m$.
The rotations $R_+$ and $R_-$ have angles $\theta=2\arctan t$ and
$-\theta$, where $0<\theta<\pi/4$. Their differences from each other
and from zero are nonzero and smaller than $\pi/2$ in absolute value.
Hence the three squares are disjoint. They lie on the unit circle,
and $Q_0$ already affinely spans the plane.

The coordinate field of $X$ is $F=\Q(t)$: its coordinates include
$u,v$, and $t=v/(1+u)$. On this field define the derivation
\[
 \partial=\frac{1+t^2}{2}\frac{d}{dt}.
\]
It satisfies $\partial u=-v$ and $\partial v=u$.
Write $R_0=I_2$, label the points by $x_{\epsilon,q}=R_\epsilon q$
for $\epsilon\in\{+,-,0\}$ and $q\in Q_0$, and assign weights
$\beta_{+,q}=\beta_{-,q}=1$, $\beta_{0,q}=-2$.
For these twelve labels put $p_i=(1,x_i)^{\mathsf T}$.
We have the three augmented moment identities
\begin{equation}\label{cons:derivation-moments}
 \sum_i\beta_i p_ip_i^{\mathsf T}=0,\qquad
 \sum_i\beta_i(\partial p_i)p_i^{\mathsf T}=0,\qquad
 \sum_i\beta_i(\partial p_i)(\partial p_i)^{\mathsf T}
       =\begin{pmatrix}0&0\\0&4I_2\end{pmatrix}.
\end{equation}
Here the block sizes are $1$ and $2$. To verify them, use
$\sum_{q\in Q_0}q=0$ and $\sum_{q\in Q_0}qq^{\mathsf T}=2I_2$.
With $J=\left(\begin{smallmatrix}0&-1\\1&0\end{smallmatrix}\right)$,
the derivatives are
$\partial R_+=R_+J$, $\partial R_-=-R_-J$, and
$\partial R_0=0$. The first spatial moment is
$2(1+1-2)I_2=0$; the mixed moment is $2J-2J=0$;
and the derivative moment is $2I_2+2I_2=4I_2$.
The constant and vector entries vanish by the same square symmetries.

Define the $\Q$-linear map $\Phi:F\otimes_\Q F\to F$ by
$\Phi(a\otimes b)=(\partial a)(\partial b)$.
For every $P\in\operatorname{Mat}_3(B)$, the product rule and
\eqref{cons:derivation-moments} give
\begin{equation}\label{cons:derivation-obstruction}
 \sum_i\beta_i\Phi\bigl((p_i\otimes1)^{\mathsf T}
                         P(1\otimes p_i)\bigr)
       =4\sum_{\alpha=1}^2 m_F(P_{\alpha\alpha}).
\end{equation}
For completeness, expand each entry of $P$ into pure tensors
$a\otimes b$. In
$\partial(p_{i\alpha}a)\partial(p_{i\beta}b)$, the terms involving
$\partial a$ or $\partial b$ vanish after summing with $\beta_i$ by
the first two moment identities and the transpose of the second.
The remaining term is $ab(\partial p_{i\alpha})
(\partial p_{i\beta})$, and the third identity gives exactly the
right side of \eqref{cons:derivation-obstruction}.

A matrix satisfying the classification criterion would make the
left side zero and the right side $8$. This contradiction proves the
claim. The argument uses a derivation only on $F$, and imposes no
regularity condition on the colorings in the Ramsey definition.
\end{proof}
